\documentclass[11pt,letterpaper]{article}
\usepackage[margin=0.85in]{geometry}
\usepackage{lmodern}
\usepackage{hyperref}
\hypersetup{colorlinks=false, urlcolor=black, pdfborder=0 0 0}
\setlength{\parindent}{0pt}
\setlength{\parskip}{0.45em}
\raggedright
\pagestyle{empty}

\begin{document}

\noindent
{\bfseries Ke-Cheng Hsu (Clio)}\\
Ph.D.\ Candidate, Department of Economics, HKUST\\
(852) 6840\,9753 \quad $|$ \quad \href{mailto:khsu@connect.ust.hk}{khsu@connect.ust.hk}\\
\href{https://cliohsu1997.github.io/}{cliohsu1997.github.io} \quad $|$ \quad \href{https://github.com/cliohsu1997}{github.com/cliohsu1997}

\vspace{0.7em}
\noindent
October 23, 2026

\vspace{0.35em}
\noindent
Hiring Team\\
Accenture --- Song, Hong Kong\\
Re: Song --- AI-Native Analyst Program (campus)

\vspace{0.2em}
Dear Hiring Team,

I am applying for the Song AI-Native Analyst Program in Hong Kong. As a Ph.D.\ candidate in Economics at HKUST (expected June 2027), I study how digital markets shape what consumers buy, pay, and keep --- the same questions Song asks when it combines strategy, technology, AI, and creativity to reinvent marketing, sales, and service. I offer what the program calls for: AI-fluent analysis of real consumer data, quantitative problem-solving on large commercial datasets, and clear communication to non-specialist audiences. I am a Hong Kong permanent resident, a native Mandarin speaker, fluent in English, with basic Cantonese.

On Taobao apparel (2017--2023; $\sim$62M RMB revenue; six categories; $\sim$20--30\% return rates), I built and estimated a model of the full customer journey --- purchase, return-shipping-insurance take-up, and return --- and ranked return-policy designs by seller profit, consumer surplus, and welfare. I collected data at scale in Python, scraping product attributes and digitizing product images with AI-based OCR, and ran a reproducible, version-controlled pipeline on AWS and HPC/SLURM. The finding guides design directly: a stricter return policy curbs wasteful returns but can block valuable purchases, so the right policy depends on the category --- providing insurance to all shoppers works best in high-margin categories, while letting each seller choose works better in low-margin ones.

With Professors Kohei Kawaguchi (HKUST) and Kei Kawai (UTokyo), I analyze Japanese bookstore point-of-sale data --- tens of billions of records across 1{,}000+ stores --- showing that when titles stock out, shoppers switch to other titles or wait for restock rather than buying from other bookstores, and that the standard approach overstates seller costs by $\sim$20\%. Each week I distill statistics, regressions, and figures into an HTML report with charts and tables that collaborators read at a glance.

I bring the same clarity to teams. As Hall Tutor at UG Hall I, I mediated resident conflicts and led orientation, workshops, and career talks. As Teaching Assistant Coordinator in the Business School, I coached and evaluated a team of teaching assistants and held their materials to one standard. As a lecturer (MSc Math Camp, rated 4.53/5) and teaching assistant (average 4.72/5), I learned to adapt quickly when the audience or the task changes.

I would welcome the chance to grow as a consultant at Song, where technology and creativity drive client growth. My resume is attached. Thank you for your consideration.

\vspace{0.7em}
Sincerely,\\[1.0em]
Ke-Cheng Hsu (Clio)

\end{document}
